\documentclass[11pt]{article}
\usepackage{preamble}
\renewcommand{\answer}[1]{}

\begin{document}

\ladderheading{AP Statistics \textperiodcentered\ Unit 2 \textperiodcentered\ Topic 2.10 \textperiodcentered\ B: 2026-11-06 \textperiodcentered\ E: 2026-11-30}{Unusual Does Not Mean Impossible}

\focusbanner{TODAY'S PROBLEM}

Suppose the same 70 percent free-throw shooter takes 100 attempts under comparable conditions and makes 58. Let $X$ be the number of makes if the career-rate model still applies.\par\smallskip \textbf{Coach:} ``Anything under 60 makes is impossible for a 70 percent shooter, so 58 proves the career rate no longer applies.''\par \textbf{Analyst:} ``A result can be unusual without being impossible; compare 58 with the model's mean and standard deviation.''\par
\begin{center}
\textbf{Career-rate model}\par\smallskip
\begin{tabular}{|l|c|c|c|}
\hline
\textbf{Set} & \textbf{Attempts} & \textbf{Rate} & \textbf{Makes} \\ \hline
\textbf{Tonight} & 100 & 0.70 & 58 \\ \hline
\end{tabular}
\end{center}
\begin{firsttakebox}{FIRST TAKE --- before the video (5 min)}
Whose statement is more defensible, the coach's or the analyst's? Name one quantity you would check.
\workspace{0.9in}
\end{firsttakebox}

\begin{callout}{\IconBook\ CENTER, SPREAD, AND UNUSUAL COUNTS (keep on the page)}{calloutgreen}
For $X\sim\mathrm{Bin}(n,p)$, $\mu_X=np$ and $\sigma_X=\sqrt{np(1-p)}$; both use count units.
Standardize with $z=\dfrac{x-\mu_X}{\sigma_X}$. A normal approximation is reasonable when both
$np\geq10$ and $n(1-p)\geq10$. A count far from the mean may be unusual, but every value from 0 through
$n$ still has positive binomial probability when $0<p<1$.
\end{callout}

\newpage
\textbf{Finish after the video:}\par\smallskip

\dokbadge{1}~\textbf{(a)} State $n$ and $p$ for $X$ and write the formulas for the binomial mean and standard deviation.\par
\workspace{0.8in}

\dokbadge{2}~\textbf{(b)} Without calculating yet, interpret the model's mean and standard deviation in the context of many sets of 100 free throws.\par
\workspace{1.3in}

\dokbadge{3}~\textbf{(c)} Compute $\mu_X$, $\sigma_X$, and the $z$-score of 58. Check $np$ and $n(1-p)$ to decide whether a normal approximation is reasonable. Is 58 unusual? Adjudicate the two statements, distinguish ``unusual'' from ``impossible,'' and explain what the coach's claim could mislead a reader into believing.\par
\workspace{2.0in}

\begin{sentenceframebox}\raggedright\textbf{Frame:} The value 58 is \blank[0.8in] standard deviations \blank[0.7in] the mean, so it is \blank[0.9in].\end{sentenceframebox}

\begin{sentenceframebox}\raggedright\textbf{Frame:} The coach confuses \blank[1.4in] with \blank[1.2in], which could make a reader believe \blank[2.3in].\end{sentenceframebox}

\begin{summaryexitbox}{EXIT --- turn this in}
One thing the video changed about my first take: \hrulefill\par
\workspace{0.4in}
\end{summaryexitbox}

\end{document}
